
Die vorliegende Arbeit präsentiert die Entwicklung und Implementierung eines modernen, digitalen Systems zur Zugangs- und Auslasskontrolle an Schulen, das die Sicherheit und Effizienz im Schulalltag signifikant steigert. Neben der reinen Sicherung des Schulzugangs adressiert das System ein besonderes organisatorisches Problem: Das traditionelle, häufig noch papierbasierte Verlassen des Schulgeländes durch das Vorzeigen von Passierscheinen oder manuellen Listen führt bei einer großen Anzahl von Schülern zu umständlichen und fehleranfälligen Prozessen. Das hier vorgestellte System kombiniert passive RFID-Technologie (Radio Frequency Identification) mit lokaler, KI-gestützter Gesichtserkennung, um einen sicheren, fälschungssicheren und vollautomatisierten Zugangs- und Auslassprozess zu ermöglichen.

Im Kern besteht die Lösung aus einer Full-Stack-Architektur, die ein Node.js-Backend, eine SQLite-Datenbank und eine plattformübergreifende React-Native-App umfasst. Während RFID-Karten für eine schnelle Identifikation sorgen, dient die Gesichtserkennung als sekundärer Validierungsschritt, um die Weitergabe von Berechtigungen (Missbrauch von Karten) durch unbefugte Mitschüler zu verhindern. Ein besonderer Fokus liegt auf der Einhaltung der Datenschutz-Grundverordnung (DSGVO), die durch eine rein lokale Verarbeitung der biometrischen Daten auf schuleigenen Servern gewährleistet wird. In der praktischen Erprobung zeigte das System eine hohe Zuverlässigkeit und Benutzerfreundlichkeit, wobei insbesondere die intuitive Benutzeroberfläche zur Zuweisung granularer Berechtigungen für Lehrkräfte und das Sicherheitspersonal hervorzuheben ist. Die Ergebnisse belegen, dass die Digitalisierung dieses Prozesses die lästige Zettelwirtschaft abschafft, eine zukunftssichere Infrastruktur für Schulen schafft und sowohl eine maßgeschneiderte Sicherheitskontrolle erlaubt als auch den administrativen Aufwand massiv reduziert.